\section{拉格朗日情形的指标公式}
\label{sec:app-c-lagrangian-index-formula}

本节把 $\mathbb R^{2n}$ 等同于 $\mathbb C^n$ 以及
$\mathbb R^n\oplus(\mathbb R^n)^*$，并常把
$\mathbb R^n\subset\mathbb C^n$ 等同于
$\mathbb R^n\times0\subset\mathbb R^n\oplus(\mathbb R^n)^*$。

我们将从定理 C.3.1 中算子的拉格朗日化
\[
 \bar\partial_{J,T}=\frac\partial{\partial\tau}
 +J\frac\partial{\partial t}+T
\]
推导其指标；拉格朗日化作用在
\[
 W_\Delta^{1,p}=W_{\mathbb R^n}^{1,p}
 (\Theta;\mathbb R^{2n}\times\mathbb R^{2n})
 :=\{\zeta\in W^{1,p}(\Theta,\mathbb R^{4n})
       \mid \zeta(\theta)\in\Delta\text{ 对所有 }\theta\in\partial\Theta\}
\]
上。这里
\[
 \Theta=(-\infty,0]\times[0,1]\mathbin{\#}D_+^2
\]
是沿 $\{t=0\}\mathbin{\#}\{\theta=0\}\subset C_+$ “展开”的 $C_+$。注意
$\Theta$（相对于边界）共形等价于 $(-\infty,0]\times[0,1]$。为计算方便，以下在后者
上工作，并仍记 $\Theta=(-\infty,0]\times[0,1]$。

我们利用\personRobbin{}与\personSalamon{}在 (RS93) 中定义的拉格朗日道路广义马斯洛夫指标，计算这一
弗雷德霍姆算子的指标。（不经过拉格朗日化的直接证明见 (Schw00)。）

为此，回忆\cref{def:app-c-generalized-cz}中的定义
\[
 \mu_{\mathrm{CZ}}(\Psi)=\mu(\operatorname{Gr}(\Psi),\Delta),
\]
这里的辛向量空间是 $(S\times S,(-\Omega)\oplus\Omega)$。通过线性辛坐标变换，可以把
对子
\[
 (\operatorname{Gr}(\Psi),\Delta),
 \qquad(S\times S,(-\Omega)\oplus\Omega)
\]
辛地变换为
\[
 (\operatorname{Gr}(\widetilde\Psi),S\times0),
 \qquad(S\times S^*,\omega_0),
\]
其中 $\widetilde\Psi\colon S\to S^*$ 是把
$\operatorname{Gr}(\Psi)$ 看作 $T^*S\cong S\times S^*$ 的截面时相伴的线性映射。因此
\begin{equation}
\label{eq:app-c-graph-coordinate-change}
 \mu(\operatorname{Gr}(\Psi),\Delta;(-\Omega)\oplus\Omega)
 =\mu(\operatorname{Gr}(\widetilde\Psi),S\times0;\omega_0).
\end{equation}
于是把 $n$ 换成 $2n$，即可由下述一般定理得到该指标。

\begin{theorem}
\label{thm:app-c-lagrangian-half-strip-index}
设 $T,J$ 如上，并令 $\Psi_\infty$ 是方程
\begin{equation}
\label{eq:app-c-lagrangian-fundamental-solution}
 \begin{cases}
 \partial\Psi/\partial t-J(\infty,t)T(\infty,t)\Psi=0,\\
 \Psi(0)=\operatorname{id}
 \end{cases}
\end{equation}
在 $(-\infty,0]\times[0,1]$ 上的解。则
\begin{equation}
\label{eq:app-c-negative-half-strip-index}
 \operatorname{Index}\bar\partial_{J,T}
 =-\mu(\Psi\mathbb R^n,\mathbb R^n)+\frac n2,
\end{equation}
其中 $\Psi\mathbb R^n$ 是由
$t\mapsto\Psi(t)(\mathbb R^n\times0)\subset\mathbb C^n$ 给出的拉格朗日道路。
\end{theorem}

本节余下部分证明\cref{thm:app-c-lagrangian-half-strip-index}，紧密沿用 (Oh99) 附录
中的证明。与那里一样，我们不在 $(-\infty,0]\times[0,1]$ 上，而在半带
$[0,\infty)\times[0,1]$ 上考虑指标问题，并证明
\begin{equation}
\label{eq:app-c-positive-half-strip-index}
 \operatorname{Index}\bar\partial_{J,T}
 =\mu(\Psi\mathbb R^n,\mathbb R^n)+\frac n2.
\end{equation}
于是由反转道路 $\Psi$ 方向的双全纯映射
\[
 (\tau,t)\in(-\tau,1-t)
\]
立刻得到\eqref{eq:app-c-negative-half-strip-index}。
\footnote{译者注：原书在这里用成员符号 $\in$ 连接坐标对；按句意应为映射符号，
译文照录。}
从现在起记 $\Theta=[0,\infty)\times[0,1]$。

现在把作用在 $W_{\mathbb R^n}^{1,p}(\Theta;\mathbb R^{2n})$ 上的算子
$\bar\partial_{J,T}=\partial/\partial\tau+J\,\partial/\partial t+T$ 变换为定义在空间
\[
 W_{\widetilde\Lambda}^{1,p}
 :=\{\zeta\in W^{1,p}(\Theta,\mathbb R^{2n})
       \mid\zeta(\theta)\in\widetilde\Lambda(\theta)\}
\]
上的另一个算子，其中
$\widetilde\Lambda\colon\partial\Theta\to\widetilde\Lambda(n)$ 定义为
\[
 \widetilde\Lambda(\theta)=
 \begin{cases}
 \mathbb R^n,&\theta\in\partial_0\Theta,\\
 \mathbb R^n,&\theta=(\tau,0),\quad\tau\leq R+1,\\
 \widetilde\Lambda_1(\theta),&\theta=(\tau,1),\quad R+1\leq\tau\leq R+2,\\
 \Psi(1)\cdot\mathbb R^n,&\theta=(\tau,1),\quad\tau\geq R+2.
 \end{cases}
\]
这里 $\partial_0\Theta$ 是图 C.1 所示 $\Theta$ 的紧部分边界在 $\partial\Theta$ 中的
部分，而
$\widetilde\Lambda_1\colon[R+1,R+2]\to\widetilde\Lambda(n)$ 是道路
\footnote{译者注：原书把 $\widetilde\Lambda_1$ 的定义域写成 $[R+1,R+2]$，
紧接的公式却量化 $\tau\geq R-1$；译文照录。}
\[
 \widetilde\Lambda_1(\tau,1)
 :=\Psi(\rho(\tau-R-1))\cdot\mathbb R^n,
 \qquad\tau\geq R-1.
\]

\begin{figure}[htbp]
\centering
\begin{tikzpicture}[x=0.95cm,y=0.8cm]
  \fill[pattern=north east lines,pattern color=gray] (0,0) rectangle (6,1);
  \draw[->] (-0.2,0) -- (6.55,0) node[right] {$\tau$};
  \draw[->] (0,-0.35) -- (0,1.7) node[above] {$t$};
  \draw (0,1) -- (6,1);
  \draw[dashed] (5,0) -- (5,1);
  \node[left] at (0,0) {$0$};
  \node[left] at (0,1) {$1$};
  \node[below] at (5,0) {$R$};
\end{tikzpicture}
\caption{定义域 $\Theta_{0,1}$。}
\label{fig:app-c-domain-theta}
\end{figure}

把定义延拓到整个 $\Theta$，令
\[
 \widetilde\Lambda(\theta)\equiv\mathbb R^n,
 \qquad\tau\leq R-1.
\]
这里 $\rho\colon\mathbb R\to[0,1]$ 是此前使用的截断函数。

注意，由于假设 $T\equiv0$ 于 $\tau\leq R$，映射
$\Psi\colon[R,\infty)\times[0,1]\to\operatorname{Sp}(2n)$ 可定义为
\[
 \Psi(\tau,t):=\Psi(\rho(\tau-R-1)t),
 \qquad\tau\geq R,
\]
并可令 $\Psi\equiv\operatorname{id}$ 于 $\tau\leq R$，从而光滑延拓到整个 $\Theta$。
因此可以定义推前算子
\[
 \Psi_*(\bar\partial_{J,T})
 =\Psi\circ\bar\partial_{J,T}\circ\Psi^{-1}
 \colon W_{\widetilde\Lambda}^{1,p}\longrightarrow L^p.
\]
该推前具有形式
\[
 \bar\partial_{\widetilde J,\widetilde T,\widetilde\Lambda}
 =\frac\partial{\partial\tau}
 +\widetilde J\frac\partial{\partial t}+\widetilde T
 \colon W_{\widetilde\Lambda}^{1,p}\longrightarrow L^p,
\]
并满足
\[
 \widetilde T\equiv0\quad\text{若 }\tau\geq R+1\text{ 或 }\tau\leq R,
 \qquad
 \widetilde J\equiv i\quad\text{若 }\tau\leq R.
\]
显然，$\bar\partial_{J,T}$ 与
$\bar\partial_{\widetilde J,\widetilde T,\widetilde\Lambda}$ 的弗雷德霍姆指标相同，故只须
计算后者的指标。回忆我们施加了横截性
$\Psi(1)\cdot\mathbb R^n\pitchfork\mathbb R^n$，因而
\[
 \widetilde\Lambda_1(\theta)\pitchfork\mathbb R^n,
 \qquad\theta=(\tau,1),\quad\tau\geq R+2.
\]
若记
\[
 \infty_+=\lim_{\tau\to\infty}(\tau,1),
 \qquad
 \infty_-=\lim_{\tau\to\infty}(\tau,0),
\]
则映射 $\widetilde\Lambda\colon\partial\Theta\to\widetilde\Lambda(n)$ 满足
\[
 \widetilde\Lambda(\infty_-)
 =\widetilde\Lambda(\tau,R+1)=\mathbb R^n,
\]
\[
 \widetilde\Lambda(\infty_+)
 =\Psi(1)\cdot\mathbb R^n
 =\widetilde\Lambda_1(R+1),
\]
以及
\[
 \widetilde\Lambda(\tau,1)
 =\Psi(\rho(\tau-R-1))\cdot\mathbb R^n,
 \qquad R+1\leq\tau.
\]
\footnote{译者注：第一式把边界值写成
$\widetilde\Lambda(\tau,R+1)$，但 $R+1$ 不在第二坐标区间 $[0,1]$ 内；译文照录。}
依构造，道路 $t\mapsto\Psi(t)\cdot\mathbb R^n$ 与
$\tau\mapsto\widetilde\Lambda_1$ 端点固定同伦，因而有相同的马斯洛夫指标：
\begin{equation}
\label{eq:app-c-boundary-path-maslov}
 \mu(\widetilde\Lambda_1,\mathbb R^n)
 =\mu(\Psi\cdot\mathbb R^n,\mathbb R^n).
\end{equation}

利用马斯洛夫指标的分层同伦不变性，即 [RS93] 的定理 2.4，可在不改变弗雷德霍姆
性质且不改变 $\widetilde\Lambda$ 的马斯洛夫指标的情况下，把算子
$\bar\partial_{\widetilde J,\widetilde T,\widetilde\Lambda}$ 形变为
\[
 \bar\partial_{i,0,\Lambda}
 =\frac\partial{\partial\tau}+i\frac\partial{\partial t}
 =\bar\partial\colon W_\Lambda^{1,p}\longrightarrow L^p,
\]
其中 $\Lambda\colon\partial\Theta_{0,1}\to\Lambda(n)$ 定义为
\[
 \Lambda(\theta)=
 \begin{cases}
 \mathbb R^n,&\theta=(\tau,0),\\
 D(t)\cdot\mathbb R^n,&\theta=(t,0),\\
 i\cdot\mathbb R^n,&\theta=(\tau,1),
 \end{cases}
\]
并且
\[
 D(t)=\operatorname{diag}\bigl(
 e^{-(\ell+\frac12)\pi it},
 e^{-\frac12\pi it},\ldots,e^{-\frac12\pi it}\bigr)
\]
对某个整数 $\ell$ 成立。
\footnote{译者注：原书给出的前两条边界分别为 $\theta=(\tau,0)$ 与
$\theta=(t,0)$，因而在字面上都落在下边界；图 C.2 则显示第二条应描述左侧竖边。
译文照录公式。}
这里把 $\Theta_{0,1}$ 等同于图 C.2 所示的半带。
